% In this file you should put the actual content of the blueprint.
% It will be used both by the web and the print version.
% It should *not* include the \begin{document}

\part{A tight lower bound for bandits with expert advice}
\label{part:main}

\textbf{Overview}

This blueprint guides the Lean formalization of the paper \emph{A Tight Lower Bound for
Non-stochastic Multi-armed Bandits with Expert Advice} by Zachary Chase, Shinji Ito and Idan
Mehalel (COLT 2026, \cite{chase2026tight}). In the \emph{bandit with expert advice} (BwE), at each
round $t$ each of $N$ experts advises one of $K$ arms, the learner pulls an arm $I_t$ and the
adversary simultaneously assigns losses $\ell_t(r) \in \{0, 1\}$ to the arms; the learner only
observes $\ell_t(I_t)$, and the adversary may adapt to the past arms of the learner. The paper
proves that for $N \ge 2K$ the minimax pseudo-regret against the best expert is
$\Omega(\sqrt{T K \log(N/K)})$, matching the upper bound of \cite{kale2014multiarmed}
(Theorem~\ref{thm:main}). The proof reduces the identification of a \emph{special batch} of experts
(SBI) to regret minimization (Lemma~\ref{lem:reduction}), bounds a Kullback--Leibler divergence
between a mixture of product laws and a product law (Lemma~\ref{lem:kl_bound}, from
\cite{ito2024minimax}), deduces a lower bound on the number of rounds of a good SBI algorithm with
one batch (Lemma~\ref{lem:two_batch}) and then with $k$ batches (Lemma~\ref{lem:many_batches}).

\textbf{Formalization.} The protocol is that of the \emph{Lean Machine Learning} library (LML):
learners are LML algorithms whose observation is the advice, whose action is the arm and whose
feedback is the loss of the pulled arm, and adversaries are LML environments whose feedback is the
loss vector; SBI algorithms are LML identification algorithms. The statements were checked against
the paper; two of them are corrected (Remarks~\ref{rem:reduction_t_pos} and~\ref{rem:main_k}), and
a gap of the proof of Theorem~\ref{thm:main} (the case of few experts per batch) is closed by a
stronger form of Lemmas~\ref{lem:two_batch} and~\ref{lem:many_batches}
(Lemmas~\ref{lem:two_batch_general} and~\ref{lem:many_batches_general}). The lower bounds on the
number of rounds are proved through a representation of the runs by i.i.d.\ \emph{pull records}
(Lemma~\ref{lem:record_representation}), which also replaces the simulation argument of the
paper's proof of Lemma~\ref{lem:many_batches}.

\chapter{The setting (Sections 1--3)}
\label{chap:setting}

\section{Bandits with expert advice}

\begin{definition}[Protocol]
  \label{def:protocol}
  \lean{Learning.Algorithm.comapBanditFeedback, Learning.Environment.FeedbackIgnoresAction,
    Learning.Environment.FeedbackIn}
  \leanok
  Let $\iota$ be a finite set of experts and $\alpha$ a finite set of arms. A \emph{learner} is an
  algorithm whose observation in round $t$ is the advice $e_t : \iota \to \alpha$, whose action is
  the arm $I_t \in \alpha$ and whose feedback is the loss $\ell_t(I_t)$ of the pulled arm. An
  \emph{(adaptive) adversary} is an environment whose observation is the advice and whose feedback
  is the loss vector $\ell_t : \alpha \to \R$; the learner observes only the loss of the pulled arm
  (it runs as \texttt{alg.comapBanditFeedback}). The adversary sees the past arms of the learner,
  but the losses of round $t$ do not depend on the arm of round $t$
  (\texttt{FeedbackIgnoresAction}). Its losses are in $\{0,1\}$ if the loss vectors are almost
  surely in $\{0,1\}^\alpha$ (\texttt{FeedbackIn}).
\end{definition}

\begin{definition}[Pseudo-regret]
  \label{def:pseudo_regret}
  \lean{Bandits.expertPseudoRegret, Bandits.diracAdvice}
  \leanok
  \uses{def:protocol}
  The pseudo-regret of a run after $T$ rounds is
  $R_T = \max_{j \in \iota} \E\big[\sum_{t < T} \ell_t(I_t) - \ell_t(e_t(j))\big]$, written in Lean
  with the Dirac advice vectors of the deterministic advice.
\end{definition}

\section{The reduced setting (Section 2)}

\begin{definition}[Arms, experts and batches]
  \label{def:arm}
  \lean{Chase2026Tight.Arm, Chase2026Tight.Expert, Chase2026Tight.Batch,
    Chase2026Tight.specialBatch}
  \leanok
  Let $k, n \in \N$. The arms are $\Arm_k = \{0\} \cup ([k] \times \{0, 1\})$, the experts
  $\Exp_{k,n} = \{0\} \cup ([k] \times [n])$, the batches $\Bat_k = \{0\} \cup [k]$; the arms
  $(u, b)$ and the experts $(u, v)$ belong to batch $u$, arm $0$ and expert $0$ to batch $0$. The
  \emph{special batch} of $\theta \in \Exp_{k,n}$ is its batch.
\end{definition}

\begin{definition}[Advice process]
  \label{def:advice}
  \lean{Chase2026Tight.adviceOf, Chase2026Tight.uniformBits, Chase2026Tight.nextAdvice,
    Chase2026Tight.adviceKernel, Chase2026Tight.initAdvice}
  \leanok
  \uses{def:arm}
  The advice is given by bits $a^{(u,v)} \in \{0,1\}$: expert $(u, v)$ advises arm
  $(u, a^{(u,v)})$, expert $0$ advises arm $0$. The bits are initially uniform. After a round in
  which an arm of batch $u$ is pulled, the bits of batch $u$ are redrawn uniformly; the bits of
  the other batches are kept.
\end{definition}

\begin{definition}[Losses and strategies]
  \label{def:strategy}
  \lean{Chase2026Tight.bern, Chase2026Tight.lossRandomness, Chase2026Tight.correctArm,
    Chase2026Tight.lossOf, Chase2026Tight.lossLaw, Chase2026Tight.lossKernel,
    Chase2026Tight.strategy, Chase2026Tight.strategy_feedbackIgnoresAction,
    Chase2026Tight.roundLaw}
  \leanok
  \uses{def:advice, def:protocol}
  Let $0 < \eps \le 0.1$. In every round, $\ell_t(0) \sim \Ber(1/2 - \eps/2)$ and in every batch
  $u$ exactly one of the two arms, the \emph{correct arm}, has loss $0$. Under the strategy $S_0$
  the correct arm of every batch is uniform; under $S_{(u^\star, v^\star)}$ the arm advised by
  $(u^\star, v^\star)$ is correct with probability $1/2 + \eps$, and the other batches are as
  under $S_0$. All these variables are independent of each other and of the past. The adversary
  $S_\theta$ (\texttt{strategy ε θ}) has the advice process of Definition~\ref{def:advice} as
  observations and these loss vectors as feedbacks; it does not use the arm of the current round.
  $P_\theta$ (\texttt{roundLaw ε θ}) is the law of the advice and losses of a round with fresh
  advice.
\end{definition}

\section{Special batch identification (Section 3)}

\begin{definition}[SBI algorithms]
  \label{def:sbi}
  \lean{Chase2026Tight.SBIAlg, Chase2026Tight.SBIAlg.toIdentAlg, Learning.IdentAlg,
    Learning.IdentAlg.stoppingTime, Learning.IdentAlg.IsRun}
  \leanok
  \uses{def:strategy}
  An algorithm for \emph{special batch identification} (SBI) is an identification algorithm in the
  setting of the bandit with expert advice: a sampling rule (a learner), a stopping rule and an
  output rule with values in the batches. It stops at a random round and outputs a batch.
\end{definition}

\begin{definition}[Good SBI algorithms, expected number of rounds]
  \label{def:is_good}
  \lean{Chase2026Tight.IsGood, Chase2026Tight.expectedRounds, Learning.IdentAlg.IsPAC}
  \leanok
  \uses{def:sbi}
  An SBI algorithm $A'$ is \emph{good} if, under every strategy $S_\theta$ of the pool, it outputs
  the special batch of $\theta$ with probability at least $0.95$. $T(A', S)$ is the expected number
  of rounds played by $A'$ before stopping when the adversary uses $S$ (infinite if $A'$ does not
  stop almost surely).
\end{definition}

\chapter{Reduction from SBI to the bandit with expert advice (Section 3)}
\label{chap:reduction}

Throughout, $0 < \eps \le 0.1$, and $u^\star$ and $j^\star$ denote the special batch and the
\emph{special expert} of the strategy $S_\theta$: $j^\star = \theta$ (expert $0$ for $S_0$).
$N_u(T)$ is the number of rounds $t < T$ in which an arm of batch $u$ is pulled.

The lemmas of this chapter are formalized on the general arm and expert types of
Chapter~\ref{chap:main} (\texttt{strategyOn}, Definition~\ref{def:embedding}): an arm $a$ is in
the batch of its reduction $\sigma(a)$, and the special expert is $\tau'(\theta)$.

\begin{lemma}[Conditional law of the losses]
  \label{lem:loss_law_cond}
  \lean{Chase2026Tight.hasCondDistrib_feedback_strategyOn, Chase2026Tight.lossKernelOn}
  \leanok
  \uses{def:strategy, lem:feedback_ignores_action_cond}
  In a run of a learner against $S_\theta$, the loss vector of round $t$ has conditional law
  $\texttt{lossLaw ε θ}(e_t)$ given the advice $e_t$ and the arm $I_t$ of round $t$ (and also given
  the history of the previous rounds, the advice and the arm, by definition of a run).
\end{lemma}

\begin{proof}
  \leanok
  The feedback kernel of $S_\theta$ is the loss kernel applied to the advice of the round, a
  function of the advice and the arm; Lemma~\ref{lem:feedback_ignores_action_cond}.
\end{proof}

\begin{lemma}[Expected losses given the advice]
  \label{lem:mean_loss}
  \lean{Chase2026Tight.meanLoss, Chase2026Tight.meanLoss_zero,
    Chase2026Tight.meanLoss_ofBatch_of_ne, Chase2026Tight.meanLoss_mk_ofBatch,
    Chase2026Tight.meanLoss_sub_meanLoss_adviceOf_ge}
  \leanok
  \uses{def:strategy}
  Let $0 \le \eps \le 1/2$ and let the advice $e$ be given by advice bits. Under
  $\texttt{lossLaw ε θ}(e)$, the expected loss of arm $0$ is $1/2 - \eps/2$, of an arm of a batch
  other than $u^\star$ it is $1/2$, and of the arm $(u, b)$ of the special batch of $S_{(u, v)}$ it
  is $1/2 - \eps$ if $(u, v)$ advises $(u, b)$ and $1/2 + \eps$ otherwise. Hence the expected loss
  of every arm $a$ exceeds the expected loss of the arm $e(j^\star)$ advised by the special expert
  by at least $\eps/2$ if $a \notin u^\star$, and by at least $0$ otherwise.
\end{lemma}

\begin{proof}
  \leanok
  The loss of arm $0$ is a $\mathrm{Ber}(1/2 - \eps/2)$ bit; the correct arm of a batch other than
  the special batch is a uniform bit; the arm advised by $(u, v)$ is incorrect with probability
  $1/2 - \eps$ under $S_{(u,v)}$. The advice of $j^\star$ is arm $0$ for $S_0$ and an arm of batch
  $u$ for $S_{(u, v)}$; compare the five cases.
\end{proof}

\begin{lemma}[Per-round regret]
  \label{lem:per_round_regret}
  \lean{Chase2026Tight.mul_measureReal_batch_ne_le_integral, Chase2026Tight.strategyOn_obsIn}
  \leanok
  \uses{lem:loss_law_cond, lem:mean_loss}
  In a run against $S_\theta$, for every round $t$,
  $\E[\ell_t(I_t) - \ell_t(e_t(j^\star))] \ge \frac{\eps}{2} \Prob(I_t \notin u^\star)$.
\end{lemma}

\begin{proof}
  \leanok
  Condition on the advice and the arm (Lemma~\ref{lem:loss_law_cond}): the expectation is that of
  the difference of the expected losses of $I_t$ and of $e_t(j^\star)$ given the advice. The advice
  is almost surely given by advice bits (it is drawn by \texttt{adviceOf} from uniform bits), and
  Lemma~\ref{lem:mean_loss} bounds this difference from below by $\frac{\eps}{2}
  \mathbf{1}\{I_t \notin u^\star\}$.
\end{proof}

\begin{lemma}[Regret and pulls of the special batch]
  \label{lem:regret_ge_pulls}
  \lean{Chase2026Tight.mul_sub_integral_pullCount_le_expertPseudoRegret,
    Chase2026Tight.integral_pullCount_history}
  \leanok
  \uses{lem:per_round_regret, def:pseudo_regret}
  In a run against $S_\theta$, $R_T \ge \frac{\eps}{2} \E[T - N_{u^\star}(T)]$.
\end{lemma}

\begin{proof}
  \leanok
  The pseudo-regret is at least the expected regret against $j^\star$ (the losses are in
  $\{0, 1\}$, hence integrable); sum Lemma~\ref{lem:per_round_regret} over $t < T$, using
  $\E[N_{u^\star}(T)] = \sum_{t < T} \Prob(I_t \in u^\star)$.
\end{proof}

\begin{definition}[The SBI algorithm of a learner]
  \label{def:sbi_of_learner}
  \lean{Chase2026Tight.sbiOfLearner, Chase2026Tight.mostPulledBatch}
  \leanok
  \uses{def:sbi}
  Given a learner $A$ and $T^\star \in \N$, the SBI algorithm $A'$ runs $A$ (with bandit feedback)
  for $T^\star$ rounds, then outputs a most pulled batch, an element of
  $\argmax_{u} N_u(T^\star)$ (a fixed choice among the maximizers).
\end{definition}

\begin{lemma}[Runs of the SBI algorithm of a learner]
  \label{lem:sbi_of_learner_run}
  \lean{Chase2026Tight.stoppingTime_sbiOfLearner, Chase2026Tight.isRun_sbiOfLearner,
    Chase2026Tight.mostPulledBatch_eq}
  \leanok
  \uses{def:sbi_of_learner}
  The SBI algorithm $A'$ of Definition~\ref{def:sbi_of_learner} stops after $T^\star$ rounds.
  Every run of the learner $A$, with output the most pulled batch in the first $T^\star$ rounds, is
  a run of $A'$. A batch pulled in more than half of the $T^\star$ rounds is the most pulled batch.
\end{lemma}

\begin{proof}
  \leanok
  The stopping rule is ``the history has length $T^\star$'' and the output rule is deterministic.
  Two distinct batches are pulled in disjoint sets of rounds, so if $N_u(T^\star) > T^\star/2$
  every other batch has fewer pulls than $u$.
\end{proof}

\begin{lemma}[Reduction, Lemma 3.1]
  \label{lem:reduction}
  \lean{Chase2026Tight.exists_isGood_expectedRounds_le,
    Chase2026Tight.exists_isGoodOn_lintegral_stoppingTime_le}
  \leanok
  \uses{def:is_good, def:pseudo_regret, def:sbi_of_learner, lem:sbi_of_learner_run,
    lem:regret_ge_pulls, lem:exists_run_universe}
  Let $A$ be a learner whose pseudo-regret under every strategy of the pool is at most $r(T)$ after
  $T$ rounds, for all $T$, and let $T^\star \ge 1$ with $T^\star \ge 1000\, r(T^\star)/\eps$. Then
  there is a good SBI algorithm $A'$ with $T(A', S) \le T^\star$ for every strategy $S$ of the pool.
\end{lemma}

\begin{proof}
  \leanok
  Take $A'$ of Definition~\ref{def:sbi_of_learner}: it stops at $T^\star$, and its runs are runs of
  $A$ (Lemma~\ref{lem:sbi_of_learner_run}). The law of its output under $S_\theta$ can be computed
  on any run; take a run of $A$ on a probability space in the universe of the regret hypothesis
  (Lemma~\ref{lem:exists_run_universe}). By Markov's inequality and
  Lemma~\ref{lem:regret_ge_pulls},
  $\Prob(N_{u^\star}(T^\star) \le \frac12 T^\star)
  \le \frac{2\, \E[T^\star - N_{u^\star}(T^\star)]}{T^\star}
  \le \frac{4\, r(T^\star)}{\eps T^\star} \le 0.004$, and $N_{u^\star}(T^\star) > \frac{T^\star}{2}$
  makes $u^\star$ the most pulled batch (Lemma~\ref{lem:sbi_of_learner_run}). The general form,
  on general arm and expert types, is proved in the same way; the paper's form is the case of the
  identity maps.
\end{proof}

\begin{remark}
  \label{rem:reduction_t_pos}
  The paper does not state $T^\star \ge 1$, but its proof divides by $T^\star$, and the statement
  is false for $T^\star = 0$: with $r(T) = T$ the hypotheses hold, and no good algorithm stops after
  $0$ rounds (when $k \ge 1$, the law of the output does not depend on the strategy).
\end{remark}

\chapter{The two-batch game (Section 4)}
\label{chap:two_batch}

In the two-batch game, $k = 1$: arm $0$ and the two arms of batch $1$, expert $0$ and the experts
$(1, v)$, $v \in [n]$. For a law $P$, $P^T$ is the law of $T$ independent draws from $P$, and
$P_{\mix^T} = \frac1n \sum_{v \in [n]} P_{(1,v)}^T$.

\begin{lemma}[Density of a round]
  \label{lem:roundLaw_withDensity}
  \lean{Chase2026Tight.roundLaw_mk_eq_withDensity, Chase2026Tight.roundDensity,
    Chase2026Tight.pi_roundLaw_mk_eq_withDensity}
  \leanok
  \uses{def:strategy, lem:pi_withDensity}
  Let $0 \le \eps \le 1/2$. $P_{(1,v)}$ has density $f_v(g) = 1 + (2 c_v(g) - 1)\, 2\eps$ with
  respect to $P_0$, where $c_v(g) = 1$ if the arm advised by $(1, v)$ in the round $g$ is the
  correct arm of batch $1$ (it has loss $0$) and $c_v(g) = 0$ otherwise. Hence $P_{(1,v)}^T$ has
  density $g \mapsto \prod_{t < T} f_v(g_t)$ with respect to $P_0^T$.
\end{lemma}

\begin{proof}
  \leanok
  Under both strategies the loss vector of the round is a function of the loss bit $z$ of arm $0$
  and of the correct arm $(1, c)$ of batch $1$, and given the advice $e$ (with bit $b$ for the
  expert $(1, v)$), the law of $(z, c)$ on $\{0, 1\}^2$ is $\Ber(1/2 - \eps/2) \otimes
  \mathrm{Unif}$ under $P_0$ and $\Ber(1/2 - \eps/2) \otimes \mathrm{Law}(b \oplus \Ber(1/2 -
  \eps))$ under $P_{(1,v)}$: the correct arm is the advised one with probability $1/2 + \eps$
  under $P_{(1,v)}$ and $1/2$ under $P_0$. Comparing the masses of the four points, the second
  law has density $(z, c) \mapsto 1 \pm 2\eps$ (sign $+$ iff $c = b$) with respect to the first;
  densities pass to images, and to the composition-product with the law of the advice. The
  product statement is Lemma~\ref{lem:pi_withDensity}.
\end{proof}

\begin{lemma}[Expected likelihood ratios]
  \label{lem:integral_ratio}
  \lean{Chase2026Tight.lintegral_prod_roundDensity, Chase2026Tight.lintegral_roundDensity,
    Chase2026Tight.lintegral_roundDensity_mul}
  \leanok
  \uses{lem:roundLaw_withDensity, lem:pi_withDensity}
  Let $0 \le \eps \le 1/2$. For $v, v^\star \in [n]$,
  $\int \prod_{t < T} f_v(g_t) \, dP_{(1,v^\star)}^T(g)$ is $1$ if $v \ne v^\star$ and
  $(1 + 4\eps^2)^T$ if $v = v^\star$.
\end{lemma}

\begin{proof}
  \leanok
  The integral of a product of functions of independent coordinates is the product of the
  integrals (Tonelli), so it suffices to show $\E_{P_{(1,v^\star)}}[f_v] =
  \E_{P_0}[f_v f_{v^\star}]$ is $1 + 4 \eps^2$ if $v = v^\star$ and $1$ otherwise. Given the advice,
  with bits $b, b^\star$ for $(1, v)$ and $(1, v^\star)$, the correct bit is uniform under $P_0$, so
  $\E_{P_0}[f_v f_{v^\star} \mid e]$ is $\frac{(1 + 2\eps)^2 + (1 - 2\eps)^2}{2} = 1 + 4\eps^2$ if
  $b = b^\star$ and $(1 + 2\eps)(1 - 2\eps) = 1 - 4\eps^2$ otherwise. If $v = v^\star$, $b = b^\star$
  always. If $v \ne v^\star$, flipping the uniform bit of $(1, v)$ preserves the law of the advice
  and exchanges the two cases, so the expectation is
  $\frac{(1 + 4\eps^2) + (1 - 4\eps^2)}{2} = 1$.
\end{proof}

\begin{lemma}[KL bound, Lemma 4.1]
  \label{lem:kl_bound}
  \lean{Chase2026Tight.klDiv_mixture_pi_roundLaw_le}
  \leanok
  \uses{lem:klDiv_mixture_le, lem:roundLaw_withDensity, lem:integral_ratio}
  If $0 < \eps \le 0.1$, then for all $T$ and $n \ge 1$,
  $\KL(P_{\mix^T} \,\|\, P_0^T) \le \frac{(1 + 4\eps^2)^T - 1}{n}$. (The paper assumes $T \ge 1$;
  the bound also holds for $T = 0$.)
\end{lemma}

\begin{proof}
  \leanok
  By Lemma~\ref{lem:roundLaw_withDensity}, $P_{(1,v)}^T$ has density
  $F_v(g) = \prod_{t < T} f_v(g_t)$ with respect to $P_0^T$. By Lemma~\ref{lem:klDiv_mixture_le}
  with weights $1/n$ and Lemma~\ref{lem:integral_ratio},
  $\KL(P_{\mix^T} \| P_0^T) \le \frac1{n^2} \sum_{v^\star} \sum_v \int F_v \, dP_{(1,v^\star)}^T - 1
  = \frac{n (1 + 4\eps^2)^T + n (n - 1)}{n^2} - 1 = \frac{(1 + 4\eps^2)^T - 1}{n}$. (The paper
  bounds the divergence by $\log\big(1 + \frac{(1 + 4\eps^2)^T - 1}{n}\big)$ with Jensen's
  inequality, then uses $\log(1 + x) \le x$; the bound holds for $T = 0$ as well.)
\end{proof}

\begin{definition}[Stopping at the $M$-th pull of a batch]
  \label{def:stop_at_pulls}
  \lean{Chase2026Tight.stopAtPulls}
  \leanok
  \uses{def:sbi}
  Let $A'$ be an SBI algorithm, $u \in [k]$ and $M \in \N$. The identification algorithm
  $B(A', u, M)$, with output in $\{\mathrm{true}, \mathrm{false}\}$, runs the sampling rule of
  $A'$, stops when $A'$ stops or as soon as it has pulled batch $u$ $M$ times, and outputs
  $\mathrm{true}$ iff it has pulled batch $u$ fewer than $M$ times and the output of $A'$ on its
  history (drawn from the output rule of $A'$) is batch $0$.
\end{definition}

\begin{lemma}[Output of the stopped algorithm]
  \label{lem:stop_at_pulls_output}
  \lean{Chase2026Tight.outputMeasure_stopAtPulls_le,
    Chase2026Tight.measure_le_outputMeasure_stopAtPulls_add}
  \leanok
  \uses{def:stop_at_pulls}
  Let $A'$ be an SBI algorithm, $u \in [k]$, $M \in \N$, $B = B(A', u, M)$ and $\tau$ the number of
  rounds played by $A'$. In every environment, $B$ outputs $\mathrm{true}$ with probability at most
  the probability that $A'$ outputs batch $0$. In every run of $A'$, the probability of
  $N_u(\tau) < M$ is at most the probability that $B$ outputs $\mathrm{true}$ plus the probability
  that $A'$ does not output batch $0$ (in the same environment).
\end{lemma}

\begin{proof}
  \leanok
  Along a run of the common sampling rule, $B$ stops at $\tau_B = \min(\tau, \tau_M)$, where
  $\tau_M$ is the time of the $M$-th pull of batch $u$. If the history at $\tau_B$ has fewer than
  $M$ pulls of batch $u$, then $\tau_B = \tau$ and the stopped histories of $A'$ and $B$ coincide;
  otherwise $B$ outputs $\mathrm{false}$. Hence, given the stopped histories, the probability that
  $B$ outputs $\mathrm{true}$ is at most the probability that $A'$ outputs batch $0$; integrate
  against the law of the run. If $N_u(\tau) < M$, then batch $u$ is pulled fewer than $M$ times up
  to $\tau$, so $\tau_B = \tau$ and the stopped histories coincide, with fewer than $M$ pulls of
  batch $u$: given them, the probabilities that $B$ outputs $\mathrm{true}$ and that $A'$ does not
  output batch $0$ sum to $1$. Integrate (Markov's inequality).
\end{proof}

\begin{lemma}[Choice of the number of pulls]
  \label{lem:two_batch_arith}
  \lean{Chase2026Tight.one_add_pow_floor_sub_one_div_le}
  \leanok
  Let $n \ge 1$, $0 < \eps \le 0.1$ and $M = \lfloor \ln(1 + n/8)/(4\eps^2) \rfloor + 1$. Then
  $\frac{(1 + 4\eps^2)^M - 1}{n} \le \frac14$.
\end{lemma}

\begin{proof}
  \leanok
  $(1 + 4\eps^2)^M \le e^{4\eps^2 M} \le e^{\ln(1 + n/8) + 4\eps^2} = (1 + n/8)\, e^{4\eps^2}$ and
  $e^{4\eps^2} \le 1 + 0.04 + 0.04^2 \le 1.05$, so
  $(1 + 4\eps^2)^M - 1 \le 0.05 + 1.05\, n/8 \le n/4$.
\end{proof}

\begin{lemma}[One batch of a general game]
  \label{lem:two_batch_general}
  \lean{Chase2026Tight.ofReal_le_lintegral_pullsBefore}
  \leanok
  \uses{lem:kl_bound, lem:events_before_pulls, lem:pinsker_event, def:is_good, def:record,
    def:stop_at_pulls, lem:stop_at_pulls_output, lem:two_batch_arith}
  Let $k \ge 1$, $n \ge 1$, $0 < \eps \le 0.1$, $u \in [k]$ and $A'$ a good SBI algorithm. Then
  $\E_{S_0}[N_u(\tau)] \ge \frac{\ln(1 + n/8)}{8 \eps^2}$, where $\tau$ is the number of rounds
  played by $A'$.
\end{lemma}

\begin{proof}
  \leanok
  Let $M = \lfloor \ln(1 + n/8)/(4\eps^2) \rfloor + 1 > \ln(1 + n/8)/(4\eps^2)$ and
  $B = B(A', u, M)$ (Definition~\ref{def:stop_at_pulls}), which stops at the latest at its $M$-th
  pull of batch $u$. Let $Q_\theta$ be the law of the output of $B$ under $S_\theta$ and
  $Q_\mix = \frac1n \sum_{v \in [n]} Q_{(u, v)}$. The output of $B$ is a function of its observed
  history up to its stopping time, which only uses the first $M$ records of batch $u$
  (Definition~\ref{def:record}): by Lemma~\ref{lem:events_before_pulls}, the data processing
  inequality, Lemma~\ref{lem:kl_bound} and Lemma~\ref{lem:two_batch_arith},
  $\KL(Q_\mix \| Q_0) \le \KL(P_{\mix^M} \| P_0^M) \le \frac{(1 + 4\eps^2)^M - 1}{n} \le \frac14$,
  and by Lemma~\ref{lem:pinsker_event}
  $\abs{Q_0(\mathrm{true}) - Q_\mix(\mathrm{true})} \le \sqrt{1/8} \le 0.4$. Since $A'$ is good and
  the special batch of $(u, v)$ is $u \ne 0$, Lemma~\ref{lem:stop_at_pulls_output} gives
  $Q_{(u, v)}(\mathrm{true}) \le 0.05$ for every $v$, hence $Q_\mix(\mathrm{true}) \le 0.05$ and
  $Q_0(\mathrm{true}) \le 0.45$. Since $A'$ outputs a batch other than $0$ with probability at most
  $0.05$ under $S_0$, Lemma~\ref{lem:stop_at_pulls_output} gives
  $P_0(N_u(\tau) < M) \le 0.45 + 0.05 = \frac12$. By Markov's inequality,
  $\E_0[N_u(\tau)] \ge M\, P_0(N_u(\tau) \ge M) \ge \frac M2 \ge \frac{\ln(1 + n/8)}{8\eps^2}$.
\end{proof}

\begin{lemma}[Pulls and rounds]
  \label{lem:pulls_le_rounds}
  \lean{Chase2026Tight.sum_pullsBefore_le}
  \leanok
  For every (random) time $\tau$, $\sum_{u \in [k]} N_u(\tau) \le \tau$.
\end{lemma}

\begin{proof}
  \leanok
  Each round before $\tau$ is a pull of at most one batch $u \in [k]$.
\end{proof}

\begin{lemma}[Lower bound for the two-batch game, Lemma 4.2]
  \label{lem:two_batch}
  \lean{Chase2026Tight.expectedRounds_ge_of_isGood_one}
  \leanok
  \uses{lem:two_batch_general, lem:pulls_le_rounds}
  Suppose that $k = 1$, $0 < \eps \le 0.1$ and $n \ge 10$, and let $A'$ be good. Then
  $T(A', S_0) \ge \frac{\ln(n/10)}{8\eps^2}$.
\end{lemma}

\begin{proof}
  \leanok
  The number of rounds is at least $N_1(\tau)$ (Lemma~\ref{lem:pulls_le_rounds}), and
  $\ln(n/10) \le \ln(1 + n/8)$; Lemma~\ref{lem:two_batch_general}.
\end{proof}

\begin{remark}
  The paper's proof treats the learner of the two-batch game as passive (it always pulls batch $1$,
  so the rounds are i.i.d.\ draws from $P_\theta$). A learner may also pull arm $0$, which does not
  redraw the advice of batch $1$: it can look at the advice of its next pull of batch $1$ before
  deciding to stop. The record representation (Lemma~\ref{lem:record_representation}) handles
  this: the advice of the next pull is uniform under every strategy.
\end{remark}

\chapter{The general SBI game (Section 5)}
\label{chap:many_batches}

\begin{lemma}[General lower bound, all $n$]
  \label{lem:many_batches_general}
  \lean{Chase2026Tight.ofReal_le_lintegral_stoppingTime}
  \leanok
  \uses{lem:two_batch_general, lem:pulls_le_rounds}
  Let $n \ge 1$, $0 < \eps \le 0.1$ and $A'$ a good SBI algorithm. Then
  $T(A', S_0) \ge k \frac{\ln(1 + n/8)}{8\eps^2}$.
\end{lemma}

\begin{proof}
  \leanok
  The number of rounds $\tau$ is at least $\sum_{u \in [k]} N_u(\tau)$
  (Lemma~\ref{lem:pulls_le_rounds}); sum Lemma~\ref{lem:two_batch_general} over the $k$ batches.
\end{proof}

\begin{lemma}[Lower bound for the general SBI game, Lemma 5.1]
  \label{lem:many_batches}
  \lean{Chase2026Tight.expectedRounds_ge_of_isGood}
  \leanok
  \uses{lem:many_batches_general}
  Let $0 < \eps \le 0.1$, $n \ge 10$ and $A'$ a good SBI algorithm. Then
  $T(A', S_0) \ge k \frac{\ln(n/10)}{20\eps^2}$.
\end{lemma}

\begin{proof}
  \leanok
  $\ln(n/10) \le \ln(1 + n/8)$ and $1/20 \le 1/8$; Lemma~\ref{lem:many_batches_general}.
\end{proof}

\begin{remark}
  The paper states Lemma~\ref{lem:many_batches} without $0 < \eps \le 0.1$ and $n \ge 10$, standing
  assumptions of its Sections 2 and 4. Its proof builds from $A'$ a good algorithm $B'$ for the
  two-batch game that simulates the other batches, and applies Lemma~\ref{lem:two_batch} to $B'$.
  Here the same conclusion is reached by applying the record argument of
  Lemma~\ref{lem:two_batch_general} to each batch of the general game, which needs no
  simulation of algorithms.
\end{remark}

\chapter{The lower bound for bandits with expert advice (Section 6)}
\label{chap:main}

\begin{definition}[Embedding of the reduced setting]
  \label{def:embedding}
  \uses{def:strategy}
  Let $K \ge 3$ and $N \ge 2K$, $k = \lfloor (K - 1)/2 \rfloor \ge 1$ and
  $n = \lfloor (N - 1)/k \rfloor \ge 1$. Embed $\Arm_k$ into $[K]$ and $\Exp_{k,n}$ into $[N]$. The
  strategy $S_\theta$ on $[K]$ and $[N]$ gives loss $1$ to the extra arms and the advice of
  expert $0$ to the extra experts.
\end{definition}

\begin{lemma}[Reduction in the embedded setting]
  \label{lem:reduction_ext}
  \uses{def:embedding, lem:reduction, lem:many_batches_general}
  Lemmas~\ref{lem:reduction} and~\ref{lem:many_batches_general} hold for learners and SBI
  algorithms on $[K]$ and $[N]$ against the embedded strategies.
\end{lemma}

\begin{proof}
  An extra arm has loss $1$, so its per-round regret against the special expert is at least
  $\eps/2$ (Lemma~\ref{lem:per_round_regret}); pulling it carries no information on the special
  batch, so the record representation (Lemma~\ref{lem:record_representation}) is unchanged.
\end{proof}

\begin{theorem}[Main lower bound, Theorem 6.1]
  \label{thm:main}
  \lean{Chase2026Tight.exists_pseudoRegret_ge}
  \leanok
  \uses{lem:reduction_ext}
  There is a universal constant $c > 0$ such that for all $K \ge 3$, $N \ge 2K$ and
  $T \ge K \ln(N/K)$, and every learner for the bandit with expert advice, there is an adaptive
  adversary with losses in $\{0, 1\}$ under which the pseudo-regret of the learner after $T$ rounds
  is at least $c \sqrt{T K \log(N/K)}$.
\end{theorem}

\begin{proof}
  Let $k, n$ as in Definition~\ref{def:embedding}, $L = \ln(1 + n/8)$ and
  $\eps = \sqrt{k L / (50 T)}$. Since $n \le 4N/K$, $L \le \ln(1 + N/(2K)) \le \ln(N/K)$ and
  $k L \le \frac{K}{2} \ln(N/K) \le T/2$, so $\eps \le 0.1$. Suppose that the pseudo-regret of the
  learner is at most $r(T) = \sqrt{T k L}/C$ under every strategy of the pool. Then
  $1000\, r(T)/\eps \le T$ for $C$ large, and Lemma~\ref{lem:reduction_ext} gives a good SBI
  algorithm with at most $T$ rounds in expectation, while Lemma~\ref{lem:many_batches_general} gives
  at least $k L/(8\eps^2) = 50T/8 > T$: a contradiction. So some strategy gives pseudo-regret
  $R_T > \sqrt{T k L}/C$. Finally $k \ge K/4$, and $L \ge \frac14 \ln(N/K)$ since $n \ge N/K \ge 2$.
\end{proof}

\begin{remark}
  \label{rem:main_k}
  The paper states Theorem~\ref{thm:main} for positive $K$. It is false for $K = 1$ (with one arm
  the pseudo-regret of every learner is $0$), and the proof needs $K \ge 3$ (one batch at least).
  The paper also assumes without loss of generality that $n > 10$, which fails when $N/K$ is
  small; Lemma~\ref{lem:many_batches_general}, valid for every $n \ge 1$, removes this assumption.
\end{remark}


\part{Prerequisites to be added to Mathlib and LML}
\label{part:prereq}

\chapter{Information theory}
\label{chap:prereq_information}

\begin{lemma}[KL divergence of a mixture]
  \label{lem:klDiv_mixture_le}
  Let $P_1, \dots, P_n$ and $Q$ be probability measures with $P_i \ll Q$, and $w_1, \dots, w_n \ge 0$
  weights with $\sum_i w_i = 1$. Then
  $\KL\big(\sum_i w_i P_i \,\big\|\, Q\big) \le \sum_i w_i \log\Big(\sum_j w_j \int \frac{dP_j}{dQ}
  \, dP_i\Big)$.
\end{lemma}

\begin{proof}
  With $M = \sum_j w_j P_j$, $\frac{dM}{dQ} = \sum_j w_j \frac{dP_j}{dQ}$ and
  $\KL(M \| Q) = \sum_i w_i \int \log \frac{dM}{dQ} \, dP_i$. Apply Jensen's inequality to the
  concave function $\log$ under each $P_i$.
\end{proof}

\begin{lemma}[Products of measures with densities]
  \label{lem:pi_withDensity}
  For finite measures $Q$ and measurable $f \ge 0$, the product of $T$ copies of
  $Q.\mathrm{withDensity}(f)$ is $Q^T.\mathrm{withDensity}\big(g \mapsto \prod_{t < T} f(g_t)\big)$.
\end{lemma}

\begin{proof}
  Both measures agree on boxes; uniqueness of product measures.
\end{proof}

\begin{lemma}[Pinsker's inequality for Bernoulli laws]
  \label{lem:pinsker_bernoulli}
  For $p, q \in [0, 1]$, $\kl(p, q) \ge 2 (p - q)^2$, where
  $\kl(p, q) = p \log\frac{p}{q} + (1 - p) \log \frac{1 - p}{1 - q}$ is the divergence between
  $\Ber(p)$ and $\Ber(q)$.
\end{lemma}

\begin{proof}
  For fixed $p$, $g(q) = \kl(p, q) - 2(p - q)^2$ vanishes at $q = p$ and
  $g'(q) = \frac{q - p}{q(1 - q)} - 4(q - p)$ has the sign of $q - p$, since $q(1 - q) \le 1/4$.
\end{proof}

\begin{lemma}[Pinsker's inequality for events]
  \label{lem:pinsker_event}
  \uses{lem:pinsker_bernoulli}
  For probability measures $P, Q$ and a measurable event $E$,
  $\abs{P(E) - Q(E)} \le \sqrt{\KL(P \| Q)/2}$.
\end{lemma}

\begin{proof}
  By the data processing inequality (Mathlib's \texttt{klDiv\_map\_le}) applied to the indicator of
  $E$, $\kl(P(E), Q(E)) \le \KL(P \| Q)$; Lemma~\ref{lem:pinsker_bernoulli}.
\end{proof}

\chapter{Sequential learning}
\label{chap:prereq_sequential}

\section{Adversaries that move simultaneously with the learner}

\begin{lemma}[Feedback independent of the current action]
  \label{lem:feedback_ignores_action_cond}
  \uses{def:protocol}
  If an environment satisfies \texttt{FeedbackIgnoresAction}, with feedback kernels
  $\kappa_t(h, o)$ depending on the history $h$ and the observation $o$ of round $t$, then in every
  run the feedback of round $t$ has conditional law $\kappa_t(H_t, O_t)$ given the history $H_t$,
  the observation $O_t$ and the action $X_t$ of round $t$; in particular it is conditionally
  independent of $X_t$ given $(H_t, O_t)$.
\end{lemma}

\begin{proof}
  The conditional law of the feedback given $(H_t, O_t, X_t)$ is the feedback kernel of the
  environment (LML's \texttt{IsAlgEnvSeq}), which factors through $(H_t, O_t)$.
\end{proof}

\section{Representation of the runs by pull records}

Fix $k, n$, $0 < \eps \le 0.1$ and a batch $u \in [k]$. A \emph{pull record} of batch $u$ is a
pair $(a, c)$ of advice bits $a \in \{0,1\}^{[n]}$ of batch $u$ and a correct arm $c \in \{0, 1\}$.

\begin{definition}[Law of the records]
  \label{def:record}
  \uses{def:strategy}
  $\rho^u_\theta$ is the law of a record of batch $u$ under $S_\theta$: $a$ uniform and, given $a$,
  $c$ uniform if $\theta$ is not an expert of batch $u$, and $c = a_v$ with probability
  $1/2 + \eps$ if $\theta = (u, v)$. It is the image, by the projection on the advice and the
  correct arm of the batch, of the law $P_{\theta'}$ of a round of the two-batch game, where
  $\theta' = (1, v)$ if $\theta = (u, v)$ and $\theta' = 0$ otherwise.
\end{definition}

\begin{lemma}[Pathwise realization of algorithms]
  \label{lem:algorithm_realization}
  An algorithm whose action space is finite can be run on a probability space carrying i.i.d.\
  uniform variables $(U_t)$ independent of the environment's randomness: there are measurable maps
  $\phi_t$ with $X_t = \phi_t(H_t, O_t, U_t)$, and likewise for the output rule of an
  identification algorithm with a finite output space.
\end{lemma}

\begin{proof}
  Inverse of the cumulative distribution function, for a finite space enumerated in a fixed order.
\end{proof}

\begin{lemma}[Record representation]
  \label{lem:record_representation}
  \uses{def:record, lem:algorithm_realization, lem:loss_law_cond}
  For every learner (or SBI algorithm) and every $\theta$, there is a probability space carrying
  independent families: i.i.d.\ record sequences $(R^{u'}_j)_{j \ge 1}$ of law $\rho^{u'}_\theta$ for
  every batch $u'$, i.i.d.\ losses of arm $0$, the randomness of the unobserved losses, and the
  learner's randomness, all but the records with laws not depending on $\theta$, and a run of the
  learner against $S_\theta$ whose observed history (advice, arms, observed losses, output) is a
  measurable function, not depending on $\theta$, of these variables. In the $j$-th pull of batch
  $u'$, the advice and the correct arm of batch $u'$ are those of $R^{u'}_j$. By the uniqueness of
  the law of a run (LML's \texttt{isAlgEnvSeq\_unique}), the observed history of every run against
  $S_\theta$ has the law of this representation.
\end{lemma}

\begin{proof}
  Build the advice of batch $u'$ in a round as the advice of the record of its next pull, the
  correct arm of a pulled batch from its record, the other losses from the independent randomness
  through \texttt{lossOf}, and the arms with Lemma~\ref{lem:algorithm_realization}. Since the advice
  of batch $u'$ is redrawn exactly after its pulls, the conditional laws of the observations and
  feedbacks given the past are those of $S_\theta$: the process is a run.
\end{proof}

\begin{lemma}[Events before the $T^\star$-th pull]
  \label{lem:events_before_pulls}
  \uses{lem:record_representation}
  Let $T^\star \in \N$ and $E$ an event of the observed history up to the stopping time $\tau$ of an
  SBI algorithm. For $\theta \in \{0\} \cup \{(u, v) : v \in [n]\}$, the probability of
  $E \cap \{N_u(\tau) \le T^\star\}$ under $S_\theta$ is
  $(\rho^{u, T^\star}_\theta \otimes \mu)(F)$ for a measurable set $F$ and a law $\mu$ that do not
  depend on $\theta$, where $\rho^{u, T^\star}_\theta = (\rho^u_\theta)^{T^\star}$ is the law of the
  first $T^\star$ records of batch $u$.
\end{lemma}

\begin{proof}
  On $\{N_u(\tau) \le T^\star\}$ the observed history only uses the first $T^\star$ records of batch
  $u$ and the advice part of record $T^\star + 1$, which is uniform under every $\theta$ and
  independent of the first $T^\star$ records; the records of the other batches have law
  $\rho^{u'}_0$ for these $\theta$. Take $\mu$ the law of all the other variables.
\end{proof}

